\chapter{DESCRIPCIÓN DEL PROBLEMA}

La incorporación del pensamiento computacional en el ámbito escolar requiere no solo propuestas pedagógicas adecuadas, sino también mecanismos concretos que permitan llevarlas a la práctica de forma organizada y sostenible. En el caso del Desafío Bebras, esta necesidad se traduce en la existencia de una infraestructura capaz de articular procesos como la difusión, el registro de participantes, la administración de tareas, la ejecución de pruebas y la publicación de resultados. Sin una plataforma que integre estos componentes, la realización del desafío depende de soluciones fragmentadas y procedimientos manuales difíciles de sostener en el tiempo.

A escala internacional, Bebras funciona como una comunidad de organizaciones responsables de ejecutar un desafío nacional en cada país. De acuerdo con los estatutos de la comunidad, solo puede existir un \textit{National Bebras Organizer} por país, y dicho organizador debe encargarse de realizar anualmente el concurso, hacerlo visible públicamente y mantener un sitio web que articule sus actividades \parencite{bebrasstatutes2022}. Además, la comunidad celebra cada año un taller internacional cuyo resultado principal es el banco anual de tareas que luego sirve de base para las ediciones nacionales \parencite{bebrasstatutes2022}. Este modelo demuestra que Bebras no opera mediante una única plataforma centralizada para todos los países, sino a través de estructuras nacionales que deben contar con capacidad propia de organización e implementación. Esta diversidad también se refleja en los sistemas tecnológicos utilizados: una presentación oficial de 2024 reporta que 29 países operan con un CMS propio, mientras que otros utilizan plataformas compartidas o adaptadas, como Cuttle, ViLLE, Moodle o sistemas desarrollados por otras organizaciones \parencite{bebrasstats2024}.

En Bolivia, esta limitación adquiere especial importancia debido a que el año 2026 se proyecta como la primera participación del país en dicha competencia internacional. Esta condición convierte a la plataforma no en un elemento complementario, sino en un requisito operativo para viabilizar la organización de una edición nacional. A diferencia de países que ya cuentan con infraestructura consolidada y procesos establecidos, el contexto boliviano enfrenta el reto inicial de construir una base tecnológica propia que permita canalizar la convocatoria, administrar el concurso y dar continuidad a futuras ediciones. La magnitud internacional del desafío refuerza esta necesidad: según estadísticas oficiales de la comunidad Bebras, en el periodo 2023--2024 participaron 3.936.642 estudiantes de 70 países \parencite{bebrasstats2024}. Este alcance evidencia que la incorporación de Bolivia supone el ingreso a una red internacional con dinámicas organizativas ya consolidadas.

El problema central radica, por tanto, en la ausencia de una plataforma web propia para la gestión integral del Desafío Bebras Bolivia. Esta carencia impide disponer de un entorno unificado para coordinar instituciones educativas, registrar participantes, estructurar exámenes por categorías, administrar ejercicios de pensamiento computacional y controlar adecuadamente la resolución de pruebas en línea. Como consecuencia, la implementación del desafío en el país se vuelve dependiente de procesos dispersos o herramientas no diseñadas específicamente para este tipo de competencia, lo que reduce la eficiencia operativa y limita las posibilidades de crecimiento de la iniciativa.

Además de la dimensión operativa, la falta de una plataforma especializada afecta la consolidación institucional del proyecto. Sin un sistema propio, resulta más difícil ofrecer una experiencia homogénea a estudiantes y docentes, gestionar contenidos académicos de manera consistente y generar condiciones para una participación recurrente en los años posteriores. En otras palabras, el problema no consiste únicamente en la ausencia de un sitio web, sino en la falta de un soporte tecnológico que permita convertir una participación inicial en un programa educativo sostenible.

Desde una perspectiva funcional, esta situación también restringe la posibilidad de adaptar el modelo internacional Bebras al contexto boliviano. La administración de categorías, la adecuación de tareas, la organización de convocatorias y la gestión de usuarios requieren herramientas específicamente pensadas para el flujo del concurso. Cuando estas capacidades no existen dentro de un sistema unificado, las tareas de coordinación se fragmentan, aumenta la carga administrativa y se dificulta asegurar condiciones consistentes para todos los participantes.

En este sentido, la descripción del problema permite establecer con claridad que la necesidad del proyecto no surge solamente del interés por desarrollar una aplicación web, sino de la exigencia concreta de contar con una plataforma que haga posible la participación inicial de Bolivia en 2026 y siente las bases para la continuidad del Desafío Bebras en el país. La solución propuesta responde, entonces, a una necesidad educativa, organizativa y tecnológica al mismo tiempo.

\clearpage
\section{ÁRBOL DEL PROBLEMA}

El árbol del problema sintetiza de manera estructurada la situación descrita anteriormente, organizando el problema central, sus principales causas y los efectos que produce en la implementación del Desafío Bebras Bolivia. Su propósito es mostrar que la ausencia de una plataforma web propia no constituye una dificultad aislada, sino el punto de convergencia de limitaciones organizativas y tecnológicas que afectan la viabilidad, eficiencia y sostenibilidad de la iniciativa en el contexto nacional.

\begin{figure}[ht]
\centering
\resizebox{\textwidth}{!}{%
\begin{tikzpicture}[
    font=\small,
    caja/.style={
        draw,
        rectangle,
        rounded corners=4pt,
        align=center,
        minimum height=1cm,
        text width=3.4cm,
        fill=white,
        inner sep=4pt
    },
    cajaGrande/.style={
        draw,
        rectangle,
        rounded corners=4pt,
        align=center,
        minimum height=1.2cm,
        text width=6cm,
        fill=white,
        inner sep=4pt
    },
    cajaCausa/.style={
        draw,
        rectangle,
        rounded corners=4pt,
        align=center,
        minimum height=1cm,
        text width=5cm,
        fill=white,
        inner sep=4pt
    },
    linea/.style={-{Stealth[length=4mm]}, thick}
]

% Efecto superior
\node[cajaGrande] (impacto) at (0,0)
    {Impacto educativo reducido \\del programa};

% Efectos intermedios
\node[caja] (exposicion) at (-9.0,-3.5)
    {Baja exposición\\ estudiantil};
\node[caja] (docentes) at (-4.5,-3.5)
    {Docentes con acceso limitado \\a programas};
\node[caja] (interes) at (0,-3.5)
    {Interés juvenil sin\\ espacios formales};
\node[caja] (alcance) at (4.5,-3.5)
    {Alcance \\organizativo \\restringido};
\node[caja] (continuidad) at (9.0,-3.5)
    {Continuidad reducida\\ del programa};

% Problema central
\node[cajaGrande] (problema) at (0,-7.5)
    {Limitaciones en la \\implementación del Desafío \\Bebras en Bolivia para \\estudiantes escolares};

% Causas intermedias
\node[cajaCausa] (infraestructura) at (-4.5,-11.5)
    {Infraestructura tecnológica\\ no especializada};
\node[cajaCausa] (gestion) at (4.5,-11.5)
    {Gestión fragmentada\\ del concurso};

% Causas raíz
\node[cajaCausa] (presencia) at (-4.5,-14.5)
    {Pensamiento computacional\\ con presencia marginal\\ en el aula};
\node[cajaCausa] (iniciativas) at (4.5,-14.5)
    {Iniciativas previas sin\\ continuidad};

% Flechas efectos intermedios -> efecto superior
\draw[linea] (exposicion.north) -- (impacto.south west);
\draw[linea] (docentes.north) -- (impacto.south west);
\draw[linea] (interes.north) -- (impacto.south);
\draw[linea] (alcance.north) -- (impacto.south east);
\draw[linea] (continuidad.north) -- (impacto.south east);

% Flechas problema central -> efectos intermedios (CORREGIDO)
\draw[linea] (problema.north west) -- (exposicion.south);
\draw[linea] (problema.north west) -- (docentes.south);
\draw[linea] (problema.north) -- (interes.south);
\draw[linea] (problema.north east) -- (alcance.south);
\draw[linea] (problema.north east) -- (continuidad.south);

% Flechas causas intermedias -> problema central
\draw[linea] (infraestructura.north) -- (problema.south west);
\draw[linea] (gestion.north) -- (problema.south east);

% Flechas causas raíz -> causas intermedias
\draw[linea] (presencia.north) -- (infraestructura.south);
\draw[linea] (iniciativas.north) -- (gestion.south);

\end{tikzpicture}%
}
\caption{Árbol del problema para la implementación del Desafío Bebras Bolivia}
\label{fig:arbol-problemas}
\end{figure}
